\section{\texorpdfstring{$\Gamma_*$}{Gamma}-Weighted Symmetric Functions} \label{LambdaG_section}

We now generalise the construction of the ring of symmetric functions in a way that incorporates $\Gamma_*$, which will define a ring $\LambdaG$ that will play a central role in this paper. Let $Q$ be the subring of $\mathbb{Z}\Gamma_*[x]$ consisting of polynomials whose constant term is a multiple of the identity (rather than an arbitrary element of $\mathbb{Z}\Gamma_*$). Consider the $n$-th tensor power (over $\mathbb{Z}$) of $Q$, which has an action of $S_n$ by permutation of tensor factors. As shorthand, for $c \in \mathbb{Z}\Gamma_*$ we write 
\[
x_i^r(c) = 1^{\otimes (i-1)} \otimes cx^r \otimes 1^{\otimes (n-i)}.
\]
This means that $x_i^r(c) x_i^s(c^\prime) = x_i^{r+s}(cc^\prime)$, and that $x_i^r(c) + x_i^r(c^\prime) = x_i^r(c+c^\prime)$. Since $Q$ has a grading (inherited from $\mathbb{Z}\Gamma_*[x]$), there is a grading on the ring of $S_n$ invariants of $Q^{\otimes n}$. We write $\Lambda_n(\Gamma_*)$ for the $S_n$-invariants of $Q^{\otimes n}$, and $\Lambda_n^k(\Gamma_*)$ for the degree $k$ component of $\Lambda_n(\Gamma_*)$. As before, for $m>n$ there are homomorphisms $\rho_{m,n}: \Lambda_{m}^k(\Gamma_*) \to \Lambda_{n}^k(\Gamma_*)$ which evaluate the polynomial variables of all tensor factors past the $n$-th at zero. These define an inverse system, and we write
\[
\Lambda^k(\Gamma_*) = \varprojlim \Lambda_n^k(\Gamma_*).
\]
\begin{remark}
It may seem unnatural to work with the ring $Q$ rather than the full polynomial ring $\mathbb{Z}\Gamma_*[x]$. The reason we do this is that evaluating elements of $Q$ at zero yields an element of $\mathbb{Z}$ rather than $\mathbb{Z}\Gamma_*$, so the codomain of $\rho_{m,n}$ is
\[
\Lambda_n^k(\Gamma_*) \otimes \mathbb{Z}^{\otimes (m-n)} = \Lambda_n^k(\Gamma_*),
\]
rather than
\[
\Lambda_n^k(\Gamma_*) \otimes (\mathbb{Z}\Gamma_*)^{\otimes (m-n)},
\]
which, being different from $\Lambda_n^k(\Gamma_*)$, would not allow us to construct an inverse system. Later we will account for these ``missing'' constant terms (see Theorem \ref{r_gamma_hom_thm}), using the ring $\RGamma$ which is defined in Section \ref{rgamma_section}.
\end{remark}

The ring $Q$ has a basis consisting of elements of the form $c x^j$ where $c \in \Gamma_*$ and $j \in \mathbb{Z}_{\geq 0}$ (we require $c = 1$ if $j=0$). Thus the ring $Q^{\otimes n}$ has a basis consisting of pure tensors in this basis of $Q$. We refer to such a pure tensor as a \emph{$\Gamma_*$-weighted monomial} and note that the $S_n$ action sends $\Gamma_*$-weighted monomials to other $\Gamma_*$-weighted monomials.
\begin{definition}
Let $\bs \lambda$ be a multipartition indexed by $\Gamma_*$. The \emph{$\Gamma_*$-weighted monomial symmetric polynomial}, $m_{\bs \lambda}(x_1, \ldots, x_n) \in Q^{\otimes n}$ is the sum of all $\Gamma_*$-weighted monomials in $Q^{\otimes n}$ that contain $c x^j$ with $j \geq 1$ as a tensor factor exactly $m_j(\bs \lambda(c))$ times. There is no restriction on the number of times $1$ may appear as a tensor factor.
\end{definition}

It is immediate that the degree of $m_{\bs\lambda}(x_1, \ldots, x_n)$ is $|\bs\lambda|$. Since the $\Gamma_*$-weighted monomial symmetric polynomials are orbit sums for the $S_n$ action on our basis of $Q^{\otimes n}$, it follows that $m_{\bs\lambda}(x_1, \ldots, x_n)$ with $|\bs\lambda|=k$ span $\Lambda_{n}^k(\Gamma_*)$, and the nonzero ones form a basis of this space. Moreover $m_{\bs\lambda}(x_1, \ldots, x_n)$ is nonzero as soon as there are enough tensor factors to accommodate all the basis vectors prescribed by $\bs \lambda$, i.e. as soon as $n \geq l(\bs\lambda)$. Finally, we observe that
\[
\rho_{m,n}(m_{\bs\lambda}(x_1,\ldots, x_m)) = m_{\bs\lambda}(x_1,\ldots, x_n),
\]
so these elements define an element of the inverse limit $\Lambda^k(\Gamma_*)$.
\begin{definition}
The \emph{$\Gamma_*$-weighted monomial symmetric function} $m_{\bs\lambda}$ is the element of $\Lambda^k(\Gamma_*)$ defined by the sequence of elements $m_{\bs\lambda}(x_1, \dots, x_n)$ as $n$ varies \textup{(}here $k = |\bs\lambda|$\textup{)}.
\end{definition}

It now follows that $\Lambda^k(\Gamma_*)$ is free as a $\mathbb{Z}$-module with basis $m_{\bs\lambda}$ indexed by all multipartitions $\bs\lambda$ of size $k$.
\begin{definition}
The ring of $\Gamma_*$-weighted symmetric functions is
\[
\LambdaG = \bigoplus_{k=0}^{\infty} \Lambda^k(\Gamma_*).
\]
\end{definition}

Similarly to the case of ordinary symmetric functions, the fact that evaluating a polynomial variable at zero is a ring homomorphism implies that $\LambdaG$ inherits the structure of a graded ring. In fact, if $\Gamma$ is the trivial group, then $\LambdaG = \Lambda$. Just as for $\Lambda$, formal properties of inverse limits automatically give us ring homomorphisms
\[
\LambdaG \to \Lambda_n(\Gamma_*).
\]

\begin{example} \label{degree_2_example}
We demonstrate how to multiply two degree 1 elements of $\LambdaG$. Suppose that $c_r \in \Gamma_*$ is a conjugacy class. Then $m_{(1)_{c_r}}$ is the element of $\Lambda^1(\Gamma)$ defined by the sequence of elements $\sum_{i=1}^n x_i(c_r) \in \Lambda_n^1(\Gamma_*)$. So we must express products of such elements in terms of $\Gamma_*$-weighted monomial symmetric polynomials $m_{\bs\lambda}(x_1, \ldots, x_n)$. Recall that we place subscripts on partitions to indicate multipartitions. For any number of variables, we have
\begin{eqnarray*}
m_{(1)_{c_r}}^2 &=& \left( \sum_i x_i(c_r) \right)^2 \\
&=& \sum_i x_i^2(c_r^2) +  \sum_{i \neq j} x_i(c_r) x_j(c_r) \\
&=& \sum_i \sum_{c_s \in \Gamma_*} A_{r,r}^s x_i^2(c_s) + 2 \sum_{i<j} x_i(c_r) x_j(c_r) \\
&=& \sum_{c_s \in \Gamma_*} A_{r,r}^s m_{(2)_{c_s}} + 2 m_{(1,1)_{c_r}}
\end{eqnarray*}
(where $A_{i,j}^k$ are the structure constants of $\mathbb{Z}\Gamma_*$), and therefore the equality between the first and last quantities may be interpreted as holding in $\Lambda^2(\Gamma_*)$. Similarly, if $c_r,c_s \in \Gamma_*$ are distinct conjugacy classes, then 
\begin{eqnarray*}
m_{(1)_{c_r}}m_{(1)_{c_s}} &=& \left( \sum_i x_i(c_r) \right)\left( \sum_j x_j(c_s) \right) \\
&=& \sum_i x_i^2(c_rc_s) +  \sum_{i \neq j} x_i(c_r) x_j(c_s) \\
&=& \sum_{c_t \in \Gamma_*} A_{r,s}^t m_{(2)_{c_t}} +  m_{(1)_{c_r}(1)_{c_s}}.
\end{eqnarray*}
\end{example}

Analogously to how $\mathbb{Z}$-linear combinations of conjugacy-class sums and $\mathbb{Z}$-linear combinations of central idempotents in $\mathbb{C}\Gamma$ define different integral forms of $\mathbb{C}\Gamma_*$, the ring $\LambdaG$ is a different integral form of the $|\Gamma^*|$-th tensor power of $\Lambda$.
\begin{proposition} \label{complex_bas_change_prop}
We have
\[
\mathbb{C} \otimes \LambdaG = \mathbb{C} \otimes \Lambda^{\otimes |\Gamma^*|}.
\]
\end{proposition}
\begin{proof}
We use the fact that $\mathbb{C} \otimes \mathbb{Z}\Gamma_* = \mathbb{C} \Gamma_* = \mathbb{C}^{\Gamma^*}$, where standard basis vectors in $\mathbb{C}^{\Gamma^*}$ are the orthogonal central idempotents $e_\chi$ in the group algebra $\mathbb{C}\Gamma$ associated to irreducible representations $\chi \in \Gamma^*$. Then $e_\chi e_{\psi} = \delta_{\chi, \psi} e_{\chi}$ (where $\delta_{\chi, \psi}$ is the Kronecker delta). Working over $\mathbb{C}$ allows us to define elements $m_{\bs\lambda}^{irr} \in \LambdaG$ similar to the $\Gamma_*$-weighted monomial symmetric functions, but instead of taking sums of pure tensors in the basis $\Gamma_*$ of $\mathbb{C}\Gamma_*$, we use the basis $e_\chi$. So $x_i^r (e_\chi)$ form a basis of (the positive degree part of) $\mathbb{C} \otimes Q$, and $m_{\bs\lambda}^{irr}$ is constructed by taking a sum of a $S_n$-orbit of a pure tensor of elements of this basis. Hence $\bs\lambda \in \mathcal{P}(\Gamma^*)$ is indexed by $\Gamma^*$ rather than $\Gamma_*$. The upshot of this is that since 
\[
x_i^r(e_\chi) x_i^s( e_\psi) = \delta_{\chi, \psi} x_i^{r+s}(e_\chi),
\]
terms corresponding to different $\chi \in \Gamma^*$ do not interact. More precisely, suppose that $\bs \lambda$ is a multipartition and let ${\bs \lambda}_\chi = {\bs \lambda}(\chi)_{\chi}$ be the multipartition concentrated in type $\chi$ taking the value $\bs \lambda(\chi)$. Then we have
\[
m_{\bs\lambda}^{irr} = \prod_{\chi \in \Gamma^*} m_{{\bs\lambda}_\chi}^{irr}.
\]
This is because each monomial in $m_{\bs\lambda}^{irr}$ arises uniquely as a product of monomials from each $m_{{\bs\lambda}_\chi}^{irr}$, each involving distinct variables. Conversely, any product of monomials from $m_{{\bs\lambda}_\chi}^{irr}$ involving the same variable for different $\chi$ is zero because $x_i^r(e_\chi) x_i^s( e_\psi) = \delta_{\chi, \psi} x_i^{r+s}(e_\chi)$. This immediately shows that
\[
\mathbb{C} \otimes \LambdaG = \bigotimes_{\chi \in \Gamma^*} \Lambda(\chi),
\]
where $\Lambda(\chi)$ has $\mathbb{C}$-basis $m_{\bs \lambda}^{irr}$, where $\bs\lambda$ is concentrated in type $\chi$. Finally, it remains to observe that $\Lambda(\chi)$ is a ring isomorphic to $\mathbb{C} \otimes \Lambda$ via the $\mathbb{C}$-linear map taking $m_{\bs\lambda}^{irr}$ to $m_{\bs\lambda(\chi)}$. This is clearly a bijection. To see that it respects multiplication note that the multiplication in $\Lambda(\chi)$ is determined by the following rule for products in $\mathbb{C} \otimes Q$: $x_i^r(e_\chi) x_i^s( e_\chi) = x_i^{r+s}(e_\chi)$. In $\Lambda$ the corresponding relation reads $x_i^r \cdot x_i^s = x_i^{r+s}$.
\end{proof}

\begin{example}
Suppose that $\Gamma = C_2 = \{1, \gamma\}$ is the cyclic group of order two (so $\gamma^2=1$). We show that $\Lambda({C_2}_*)$ is not isomorphic to $\Lambda \otimes \Lambda$ as a graded ring by considering the module of indecomposables. Suppose that $A = \bigoplus_{i=0}^\infty A_i$ is a graded algebra and $I = \bigoplus_{i=1}^\infty A_i$ is the corresponding augmentation ideal. The \emph{module of indecomposables} is defined to be $I/I^2$, which is a (graded) module for $A/I$. We compute the degree two component of $I/I^2$ when $A$ is either $\Lambda({C_2}_*)$ or $\Lambda \otimes \Lambda$.

Recall that $\Lambda = \mathbb{Z}[e_1, e_2, \ldots]$. If $A = \Lambda \otimes \Lambda$, then the degree $1$ component of $A$ has basis $e_1 \otimes 1, 1 \otimes e_1$, while the degree $2$ component has basis $e_2 \otimes 1, e_1^2 \otimes 1, e_1 \otimes e_1, 1 \otimes e_2, 1 \otimes e_1^2$. Then the degree 2 component of $I^2$ has basis $e_1^2 \otimes 1, e_1 \otimes e_1, 1 \otimes e_1^2$. We conclude that the degree $2$ component of the module of indecomposables of $A$ is $\mathbb{Z}^2$ (spanned by the classes of $e_2 \otimes 1$ and $1 \otimes e_2$).

Analogously, for $A = \Lambda({C_2}_*)$, the degree $1$ component has basis $m_{(1)_1}, m_{(1)_\gamma}$, while the degree $2$ component has basis $m_{(1,1)_1}, m_{(2)_1}, m_{(1)_1 (1)_\gamma}, m_{(1,1)_\gamma}, m_{(2)_\gamma}$. The degree $2$ component of $I^2$ is spanned by 
\begin{eqnarray*}
m_{(1)_1}^2 &=& m_{(2)_1} + 2m_{(1,1)_1} \\
m_{(1)_1}m_{(1)_\gamma} &=&  m_{(2)_\gamma} + m_{(1)_1 (1)_\gamma} \\
m_{(1)_\gamma}^2 &=& m_{(2)_1} + 2 m_{(1,1)_\gamma},
\end{eqnarray*}
which we computed in Example \ref{degree_2_example}. From this we see that the degree $2$ component of $I/I^2$ is isomorphic to $\mathbb{Z}^{2} \oplus \mathbb{Z}/2\mathbb{Z}$. We conclude that $\Lambda({C_2}_*) \neq \Lambda^{\otimes 2}$. We discuss the algebraic structure of $\LambdaG$ further in the appendix.
\end{example}



